\documentclass[11pt,a4paper]{article}
\usepackage[utf8]{inputenc}
\usepackage[T1]{fontenc}
\usepackage[spanish,provide=*]{babel}
\usepackage{geometry}
\usepackage{booktabs}
\usepackage{longtable}
\usepackage{xcolor}
\usepackage{listings}
\usepackage{hyperref}
\usepackage{titlesec}
\geometry{margin=2.3cm}
\definecolor{azulWCD}{RGB}{17,67,112}
\definecolor{tituloColor}{RGB}{0,55,110}
\definecolor{reglaColor}{RGB}{0,90,160}
\hypersetup{colorlinks=true,linkcolor=azulWCD,urlcolor=azulWCD}
\lstset{basicstyle=\ttfamily\small,breaklines=true,frame=single,
  columns=fullflexible,keywordstyle=\color{azulWCD}}
\titleformat{\section}{\normalfont\Large\bfseries\color{azulWCD}}
  {\thesection.}{0.6em}{}
\titleformat{\subsection}{\normalfont\large\bfseries}
  {\thesubsection.}{0.5em}{}

\begin{document}
\pagestyle{empty}

%------------------------------------------------
% PORTADA
%------------------------------------------------

\begin{titlepage}
\begin{center}

\vspace*{0.5cm}
\textcolor{reglaColor}{\rule{\linewidth}{1.2pt}}

\vspace{1.3cm}
{\Large\scshape Universidad Industrial de Santander}\\[0.2cm]
{\large Escuela de Física --- Facultad de Ciencias}\\[0.2cm]
{\large Doctorado en Física}

\vspace{2.2cm}
{\Huge\bfseries\color{tituloColor} Revisión técnica de la\\[0.2cm]
superficie óptica del PMT}

\vspace{1.8cm}
{\large\textbf{Informe técnico}}

\vspace{0.6cm}
\begin{minipage}{0.82\linewidth}
\centering
\textit{``Simulación WCD en Geant4''}
\end{minipage}

\vfill

{\large\textbf{Autor}}\\[0.15cm]
Propuesta de revisión para MEIGA\\[0.1cm]
wcd-nacl-crns-simulations

\vspace{1.2cm}
\textcolor{reglaColor}{\rule{0.5\linewidth}{0.6pt}}

\vspace{1cm}
11 de septiembre de 2026

\vspace{1cm}
\textcolor{reglaColor}{\rule{\linewidth}{1.2pt}}

\end{center}
\end{titlepage}

\pagestyle{plain}

\section{Alcance}
Se revisaron los archivos \texttt{Materials-mejorado-optica-Tyvek-RINDEX(1).cc}
y \texttt{SaltyWCD-mejorado-optica-Tyvek-RINDEX(1).cc}. El análisis cubre la
ventana de borosilicato, la interfaz medio--PMT, la fotocátodo, la eficiencia
cuántica, la reflexión y el mecanismo de registro de fotones ópticos. No se
ejecutaron simulaciones ni se modificó un repositorio externo.

\section{Conclusión ejecutiva}
La implementación actual no constituye un modelo físico completo de una
superficie de PMT. El material denominado \texttt{Pyrex} tiene índice de
refracción constante $n=1.47$, pero una longitud de absorción de
$0.0005\,\mathrm{mm}=0.5\,\mu\mathrm{m}$. En la geometría, toda la media
elipsoide sólida se construye con ese material y se declara sensible. Por
tanto, el volumen se comporta prácticamente como un absorbente negro, no como
una ventana de vidrio seguida de una fotocátodo.

No se encontró una \texttt{G4OpticalSurface} específica del PMT, ni propiedades
\texttt{EFFICIENCY} o \texttt{REFLECTIVITY} para la fotocátodo. Tampoco existe
una interfaz explícita medio detector--PMT. La salinidad sí cambia el
\texttt{RINDEX} del medio y por ello cambiaría las pérdidas de Fresnel frente a
una ventana real; sin embargo, la parametrización actual impide estudiar de
forma separada transmisión por vidrio, reflexión en fotocátodo y detección.

\section{Hallazgos}
\begin{longtable}{p{0.25\textwidth}p{0.31\textwidth}p{0.34\textwidth}}
\toprule
Elemento & Estado encontrado & Evaluación \\
\midrule
\endhead
\texttt{RINDEX} de Pyrex & Dos puntos, 2--4 eV, $n=1.47$ constante &
Útil sólo como aproximación gruesa. No reproduce dispersión espectral y no
cubre completamente la malla del agua hasta 4.136 eV. \\
\texttt{ABSLENGTH} de Pyrex & $0.5\,\mu$m constante &
Incorrecta para una ventana de borosilicato. Implementa deliberadamente un
absorbente superficial mediante una propiedad de volumen. \\
Geometría & Media elipsoide sólida de Pyrex &
No separa pared de vidrio, vacío interno y capa de fotocátodo. El espesor óptico
no corresponde al espesor de una ventana. \\
Superficie de fotocátodo & Ausente &
No hay tipo, acabado, reflectividad ni eficiencia cuántica espectral. \\
Detector sensible & \texttt{logPMT} usa \texttt{G4MPMTAction} &
Puede contar fotones que entran o pasos dentro del volumen, pero su semántica no
está disponible en los archivos revisados. Es crítico auditarlo antes de usar
\texttt{EFFICIENCY}. \\
Posición/orientación & Elipsoide truncado entre $-z$ y 0, colocado en la parte
superior del agua & Coherente con una cara curva orientada hacia abajo. El
desplazamiento de 0.1 mm evita coincidencia exacta, pero debe verificarse con
\texttt{CheckOverlaps}. \\
Dependencia con NaCl & Cuatro materiales con diferentes \texttt{RINDEX} &
La respuesta angular de Fresnel depende automáticamente del medio. La QE propia
de la fotocátodo no debe cambiar con NaCl salvo que existan medidas específicas;
lo que sí cambia es la probabilidad de llegada y transmisión en la interfaz. \\
\bottomrule
\end{longtable}

\section{Por qué no se sustituyeron números sin datos}
Cambiar $0.5\,\mu$m por una longitud arbitraria o inventar una curva de QE
produciría resultados no trazables. Para una implementación científica deben
aportarse: modelo exacto del PMT, curva de QE o eficiencia de detección frente a
longitud de onda, reflectividad de la fotocátodo y transmisión/absorción de la
ventana. Si la curva del fabricante se digitaliza desde una gráfica, debe
guardarse el CSV, la referencia, la versión y la incertidumbre de digitalización.

\section{Arquitectura de corrección propuesta}
\subsection{Geometría}
Construir tres regiones: (1) una cáscara delgada de borosilicato, obtenida como
diferencia de dos elipsoides; (2) el interior del tubo, normalmente vacío; y
(3) una fotocátodo delgada o, preferiblemente, una superficie de frontera sobre
la cara interna relevante. La fracción activa del 90\% mencionada en el código
debe implementarse geométricamente o eliminarse del comentario.

\subsection{Ventana}
La tabla del vidrio debe contener \texttt{RINDEX(E)} y
\texttt{ABSLENGTH(E)} medidos o documentados, en orden creciente de energía y
cubriendo toda la banda de fotones. La interfaz agua/solución--vidrio puede ser
\texttt{dielectric\_dielectric}, \texttt{polished}. En esta frontera Geant4
calcula reflexión y refracción de Fresnel usando los índices de ambos medios;
por ello no se debe duplicar la reflectividad de Fresnel con una tabla ad hoc.

\subsection{Fotocátodo}
La fotocátodo puede representarse como \texttt{dielectric\_metal}, acabado
\texttt{polished}, con \texttt{REFLECTIVITY(E)} y \texttt{EFFICIENCY(E)}.
En Geant4, \texttt{EFFICIENCY} es la fracción detectada de los fotones que no
son reflejados; no debe confundirse automáticamente con una QE tabulada respecto
a fotones incidentes. Si $QE_{\rm inc}(E)$ es la probabilidad por fotón incidente
y $R(E)$ la reflectividad, la conversión compatible es
\[
  \mathrm{EFFICIENCY}(E)=\frac{QE_{\rm inc}(E)}{1-R(E)},
\]
siempre que el cociente esté entre 0 y 1 y que el modelo experimental emplee
esa convención. De lo contrario debe implementarse una aceptación explícita.

\subsection{Registro de detección}
Antes de activar la superficie se debe escoger un único método:
\begin{enumerate}
  \item detectar el estado \texttt{Detection} de \texttt{G4OpBoundaryProcess}
  en una acción de pasos; o
  \item invocar el detector sensible desde el proceso de frontera, si la versión
  de Geant4 y la configuración del proyecto lo permiten; o
  \item dejar pasar el fotón al volumen sensible y aplicar allí una sola curva
  de aceptación.
\end{enumerate}
Mezclar estos métodos puede producir doble conteo; una superficie que absorbe
el fotón antes de entrar puede también volver inalcanzable a
\texttt{G4MPMTAction}.

\section{Cambios de código requeridos}
\begin{longtable}{p{0.31\textwidth}p{0.59\textwidth}}
\toprule
Archivo & Cambio \\
\midrule
\endhead
\texttt{Materials.cc/.h} & Separar propiedades del vidrio y de la fotocátodo;
declarar la superficie del PMT; validar tamaños, orden y dominio de las tablas. \\
\texttt{SaltyWCD.cc/.h} & Crear cáscara, interior y frontera; conservar la
orientación; guardar punteros físicos necesarios para
\texttt{G4LogicalBorderSurface}; comprobar solapamientos. \\
\texttt{G4MPMTAction.cc/.h} & Confirmar qué paso cuenta, si elimina el fotón y
si ya aplica QE/CE. Adaptarlo a un único mecanismo de detección. \\
\texttt{G4WCDSteppingAction.cc/.h} & Si se escoge detección de frontera, leer
\texttt{G4OpBoundaryProcessStatus::Detection}, verificar el nombre/copia del PMT
y registrar exactamente un fotoelectrón. \\
Lista de física & Confirmar \texttt{G4OpticalPhysics}, Cherenkov, absorción,
Rayleigh/Mie y \texttt{G4OpBoundaryProcess}; revisar parámetros ópticos globales. \\
Macros & Fuente monoenergética/monocromática de \texttt{opticalphoton}, barridos
de energía y ángulo, semillas, verbosidad y número de eventos. \\
Clases de datos & Revisar \texttt{OptDevice}, \texttt{Detector}, \texttt{Event}
y la salida para distinguir fotones incidentes, reflejados, absorbidos y
fotoelectrones. \\
Construcción/acciones & Verificar registro de detector, acciones por hilo y
compatibilidad multithreading. \\
\bottomrule
\end{longtable}

\section{Pruebas mínimas de aceptación}
Para cada medio (agua, 2.5\%, 5\% y 10\% NaCl), lanzar fotones ópticos
monocromáticos normales y oblicuos a la ventana. Registrar $N_{\rm inc}$,
$N_{\rm refl}$, $N_{\rm trans}$, $N_{\rm abs}$ y $N_{\rm det}$. Exigir cierre
estadístico
\[
N_{\rm inc}=N_{\rm refl}+N_{\rm trans}+N_{\rm abs}+N_{\rm det}
\]
según las fronteras definidas. Comprobar límites: QE=0, QE=1, R=0 y R=1; luego
comparar Fresnel contra cálculo analítico. Finalmente ejecutar la geometría
completa y verificar que ningún track se cuenta dos veces.

\section{Estado de los dos archivos entregados}
No se efectuó una sustitución numérica especulativa en los dos archivos fuente.
El cambio científicamente correcto no queda autocontenido en ellos: depende de
\texttt{Materials.h}, de la geometría interna del PMT y, de manera decisiva, de
\texttt{G4MPMTAction}. Los archivos actuales deben considerarse un modelo
idealizado de absorbente sensible, no una simulación validada de superficie de
PMT. El parche debe integrarse sólo después de recuperar esos archivos y la
curva óptica del PMT seleccionado.

\section{Resumen listo para commit}
\begin{lstlisting}
docs(pmt): document optical-surface audit and required detector model

- identify sub-micron Pyrex absorption length as an idealized black absorber
- document missing photocathode EFFICIENCY and REFLECTIVITY
- define window/photocathode geometry and boundary-model requirements
- explain NaCl dependence through medium-to-window Fresnel transmission
- list G4MPMTAction, stepping action, physics list and macro audits
- specify four-medium optical-photon validation matrix
\end{lstlisting}

\section{Referencias técnicas}
Geant4 Collaboration, \emph{Book for Application Developers}, sección de
procesos ópticos y fronteras: \\
\url{https://geant4.web.cern.ch/documentation/pipelines/master/bfad_html/ForApplicationDevelopers/TrackingAndPhysics/physicsProcess.html}

Geant4 Collaboration, implementación de \texttt{G4OpBoundaryProcess}: \\
\url{https://github.com/Geant4/geant4/blob/master/source/processes/optical/src/G4OpBoundaryProcess.cc}

\end{document}
